\section{What an observer learns without the private read}
\label{sec:tech:leakage}

The previous section built a private read. This one measures what happens
without it, and it is the strongest result we have: it is the quantitative
answer to why the delivery layer must be private at all, and neither reference
paper reports it, because neither implements both halves of the system.

\subsection{Why a fetch is not one bit}
\label{sec:tech:leakage:why}

\nudgeplain{} publishes the item embedding matrix $\itemembed$ in the clear.
That is the design and not an oversight: each converged row of $\itemembed$ is
opened before the next is computed, which is what makes \textsc{SetOrthogonal}
a Gram--Schmidt against \emph{public} vectors and keeps the whole
power-iteration approach affordable under secret sharing.

The consequence is the one this section measures. $\itemembed$ is a complete
latent-factor model of the catalogue, known to everyone. So an adversary who
observes \emph{which} records a user fetched does not learn $j$ isolated facts;
each fetch is a constraint on the user's private embedding $\avec{i}$, a
projection onto a basis the adversary already holds. The question is how fast
those constraints collapse.

\subsection{The experiment}
\label{sec:tech:leakage:setup}

The adversary is given $\itemembed$ and the \emph{indices} of the $j$ records
the user fetched. It is given neither the scores nor $\userembed$. It computes
$\hat{a}$, the centroid of the fetched items' unit embeddings, and we report
$\cos(\hat{a}, \avec{i})$ together with the top-20 overlap between the ranking
$\hat{a}$ induces and the user's true ranking. All 943 users, at
$\dime = 16$, $\iters = 10$, against the same factorisation the demo serves.

\paragraph{The control is the load-bearing part of the design.} A
random-direction control would flatter any attack of this shape. Popular items
sit near the top of nearly everyone's ranking, so an estimator can score well
having learned only what is \emph{public}. The honest control therefore runs
the identical estimator on the $j$ globally most-popular items, ignoring the
user entirely, and the attack may claim only what it achieves \emph{above} that
line. The random control is retained to show where the floor actually is.

\begin{figure}[h]
\centering
\includegraphics[width=\linewidth]{figures/fig_attack}
\caption{Reconstruction of a user's private embedding from observed fetches
alone, over all 943 users. \textbf{Left:} cosine against the true embedding,
with the interquartile band on the attack. \textbf{Right:} the share of the
user's \emph{next} twenty recommendations the adversary can name. Both panels
share the same controls. The gap between the attack and the popularity control
is the part that is genuinely about the individual rather than about what is
popular, and it is stable at $+0.39$ to $+0.42$ across every $j$. Regenerated
by \texttt{mingw32-make figures} from
\texttt{bench/results/leakage\_attack.jsonl}.}
\label{fig:attack}
\end{figure}

\subsection{Result}
\label{sec:tech:leakage:result}

\begin{table}[h]
\centering
\small
\caption{Recovery against the number of observed fetches. The overlap column
excludes the $j$ items already seen, so it measures prediction of the user's
\emph{next} recommendations rather than repetition of the ones handed to the
adversary.}
\label{tab:attack}
\begin{tabular}{@{}rccc@{}}
\toprule
$j$ & $\cos(\hat{a}, a)$ & popularity control & next-20 overlap \\
\midrule
1  & \textbf{0.55} & 0.26 & 38\% \\
2  & 0.66 & 0.26 & 47\% \\
5  & 0.77 & 0.39 & 53\% \\
10 & \textbf{0.84} & 0.44 & \textbf{56\%} \\
20 & 0.89 & 0.47 & 55\% \\
50 & 0.92 & 0.53 & 46\% \\
\bottomrule
\end{tabular}
\end{table}

\textbf{A single observed fetch already gives $\cos = 0.55$, twice the
popularity control. Ten give $0.84$, and are enough to predict 56\% of that
user's next twenty recommendations.} The random control sits at $0.003$.

Two caveats we raise ourselves rather than wait to be asked.

First, \emph{an individual user's cosine is not evidence}. The random control's
interquartile spread is roughly $\pm 0.18$, near $1/\sqrt{\dime}$ at
$\dime = 16$, so one user scoring $0.2$ is inside the noise of guessing. Only
population means, and the margin over the popularity control, carry a claim.

Second, the overlap column is \emph{non-monotone}: it peaks near $j = 10$ and
falls by $j = 50$. That is a property of the metric, not a degradation of the
attack. The $j$ observed items are excluded from the candidate pool, so as $j$
grows the easy head of the ranking is progressively removed and what remains is
a harder tail. The cosine rises monotonically throughout, which is the check
that the estimate itself keeps improving.

\paragraph{An honest negative result.} We also implemented a second, stronger
estimator, which uses the fact that the fetched set is the \emph{top} of the
score vector and so every fetched item outranks every unfetched one --- strictly
more information than the centroid uses. It is \emph{worse} at every $j \ge 2$
(0.81 against 0.92 at $j = 50$). We report it rather than quietly dropping it,
because it sharpens the conclusion rather than weakening it: the cheap, obvious
attack is already sufficient, so an adversary here needs no sophistication.

\subsection{What it costs to close the gap}
\label{sec:tech:leakage:cost}

\begin{figure}[h]
\centering
\includegraphics[width=\linewidth]{figures/fig_pir_cost}
\caption{The private read against the two baselines a reader would otherwise
propose, at the ML-100K operating point (1682 items, 256\,B records,
$\mathrm{db} = 11$). Both panels must be read together; showing only the left
one would be advocacy rather than evaluation.}
\label{fig:pircost}
\end{figure}

Figure~\ref{fig:pircost} prices it. Against B5, downloading the entire
catalogue and filtering at home --- which is trivially private and is the first
thing a reader proposes, since 430\,KB is not obviously unreasonable ---
DPF-PIR sends \textbf{446$\times$ fewer bytes} counting both servers.

Against B1, a cleartext lookup with no privacy at all, it costs about
\textbf{50{,}000$\times$ more server time}. And read plainly,
\textbf{DPF-PIR is the slowest of the three in server CPU}, about
$5\times$ slower than B5's local copy. Its advantage is entirely in
communication.

Setting the extra computation against the saved bytes, the private read is the
better choice on any link slower than \textbf{8.0\,Gbit/s} --- which is to say
on every real network, but we would rather state the crossover than imply there
is not one. Absolute timings on this machine vary up to $4\times$ between
back-to-back runs, so these are ratios taken within a single run, as
Section~\ref{sec:tech:bench} explains.

\subsection{Would decoys do instead? No, and the reason is interesting}
\label{sec:tech:leakage:decoys}

The obvious cheap alternative to a private read is to fetch the $\topk$ real
recommendations plus $r$ decoys and let the observer sort it out. We measured
it, because ``obviously insufficient'' is not an argument.

Decoys must be drawn \emph{popularity-weighted}, not uniformly. Real
recommendations are popularity-skewed, so uniform decoys would come mostly from
the long tail and be separable by inspection alone --- an observer would
discard the unpopular ones and run the original attack on what was left.
Measuring a defence against an adversary too naive to notice that would flatter
it.

\begin{figure}[h]
\centering
\includegraphics[width=\linewidth]{figures/fig_decoy}
\caption{The decoy defence, and the floor it cannot cross. Even 40 decoys ---
five times the real traffic, at 48\,KB per query --- leave the adversary at
$\cos = 0.71$, far above the popularity control at $0.44$.}
\label{fig:decoy}
\end{figure}

\begin{table}[h]
\centering
\small
\caption{Cost and benefit of decoys, at $\topk = 10$. Bandwidth is both
servers, from Section~\ref{sec:tech:leakage:cost}.}
\label{tab:decoy}
\begin{tabular}{@{}rrrr@{}}
\toprule
decoys $r$ & fetches & $\cos(\hat{a}, a)$ & bytes/query \\
\midrule
0  & 10 & 0.843 & 9{,}660 \\
5  & 15 & 0.828 & 14{,}490 \\
10 & 20 & 0.806 & 19{,}320 \\
40 & 50 & \textbf{0.705} & 48{,}300 \\
\bottomrule
\end{tabular}
\end{table}

\textbf{The defence does not work, and the reason is the same fact that made
the popularity control necessary.} Because decoys must be popularity-weighted
to be unremarkable, they point roughly along the popularity direction --- and
the popularity direction is exactly what the control already recovers, at
$\cos = 0.44$. So decoys drag the estimate down \emph{towards} the control's
level and no further. There is a floor built into the defence, and it sits
where an adversary who watched nobody already stands.

Quintupling the traffic buys a drop of $0.14$ and leaves the adversary far
above that floor. A private read costs about a tenth of that bandwidth and
removes the observation entirely. This is a negative result about the
alternative, and it is the strongest argument for the mechanism this project
actually builds.

\medskip

Together, the parts of this section make the argument the composition exists to
make: without a private read, ten observed fetches recover the user; the cheap
alternative to a private read does not close the gap; and a private read closes
it for a few hundred bytes and a fraction of a millisecond. \nudgeplain{} leaves
this to ``other means''. This is what those means cost, and what the plausible
substitutes do not buy.
